\textbf{1st version}

\begin{description}
\item[1.] \emph{Determine coordinates of Lima and Yogyakarta.}
  \hfill(15 points)\\
  Assume that the equator is the $z = 0$ plane and the Earth is a unit
  sphere. The Cartesian coordinates of Lima are
  \begin{align*}
  \boldsymbol{x}_1 &= (x_1, y_1, z_1)\\
    &= \bigl(\cos(-12^\circ 2')\cos(282^\circ 58'48''),\
       \cos(-12^\circ 2')\sin(282^\circ 58'48''),\
       \sin(-12^\circ 2')\bigr)\\
    &= (0.2197308726,\ -0.9530236794,\ -0.2084807188)
  \end{align*}
  The Cartesian coordinates of Yogyakarta are
  \begin{align*}
  \boldsymbol{x}_2 = (x_2, y_2, z_2)
    &= \bigl(\cos(-7^\circ 47')\cos(110^\circ 26'),\
       \cos(-7^\circ 47')\sin(110^\circ 26'),\
       \sin(-7^\circ 47')\bigr)\\
    &= (-0.3459009563,\ 0.9284459898,\ -0.1354273699)
  \end{align*}
\item[2.] \emph{Determine the normal of the plane defined by Lima,
  Yogyakarta, and the center of the Earth.} \hfill(20 points)\\
  The normal of the plane containing the two places and the center of the
  Earth is
  \[ \boldsymbol{x}_1 \times \boldsymbol{x}_2 = (a, b, c)
     = (0.3226285776,\ 0.1018712545,\ -0.12564355210) \]
\item[3.] \emph{Find the equation of the (shortest flight) path from Lima to
  Yogyakarta.} \hfill(25 points)\\
  Let the equation of the plane be $ax + by + cz = 0$ or equivalently
  \[ y = -a'x - c'z \]
  where $a' = -3.16702272082062$ and $c' = 1.23335628599724$. For the sake
  of simplicity, from now on let us denote $a'$ and $c'$ as $a$ and $c$,
  respectively. The intersection of the Earth and the plane is the great
  circle
  \[ x^2 + (-ax - cz)^2 + z^2 = 1 \quad\text{or equivalently}\quad
     (1 + a^2)x^2 + 2acxz + (1 + c^2)z^2 - 1 = 0 \]
  It is also the equation for the shortest flight path from Lima to
  Yogyakarta.
  % FIGURE NEEDED: unit sphere with Lima, Yogyakarta, the flight path and
  % the normal vector x1 x x2 (2015_th_sol.pdf, page 19).
\item[4.] \emph{Determine the $z$-value of the southernmost point of the
  path (at that point, the discriminant of the equation
  $(1 + a^2)x^2 + 2acmx + (1 + c^2)m^2 - 1 = 0$ equals zero).}
  \hfill(30 points)\\
  The plane $z = m$ intersects the great circle at either one or two
  points. The number of intersection points at the highest or at the lowest
  point is one, i.e.\ when the discriminant is zero.
  \begin{align*}
  (1 + a^2)x^2 &+ 2acmx + (1 + c^2)m^2 - 1 = 0\\
  \intertext{The discriminant is}
  D &= (2acm)^2 - 4(1 + a^2)[(1 + c^2)m^2 - 1]\\
    &= 4 + 4a^2 - (4 + 4a^2 + 4c^2)m^2
  \end{align*}
  Therefore $D = 0$ if
  $m = \pm\sqrt{\dfrac{1 + a^2}{1 + a^2 + c^2}} = 0.937444937497383$. Then
  choose
  \[ z = -0.937444937497383 \]
\item[5.] \emph{Determine the latitude of the southernmost point.}
  \hfill(10 points)\\
  Thus, the latitude of the point is
  $\sin^{-1} z = 1.21521692653748 = 69.6268010651290^\circ
   = 69^\circ 37' 36''\,S$.
\end{description}

\textbf{2nd version (using the Law of Cosines for sides)}

\begin{description}
\item[1.] \emph{Determine the lengths of the sides of the triangle}
  \hfill(20 points)\\
  \[ s_2 = 90^\circ - 12^\circ 2' = 77^\circ 58' = 1.3608\ \mathrm{rad},\quad
     s_3 = 90^\circ - 7^\circ 47' = 82^\circ 13' = 1.4350\ \mathrm{rad}; \]
  \[ A_1 = (180^\circ - 77^\circ 1') + (180^\circ - 110^\circ 26')
     = 172^\circ 33' = 3.0116\ \mathrm{rad} \]
  Using the Law of Cosines for sides
  \[ \cos s_1 = (\cos s_2)(\cos s_3) + (\sin s_2)(\sin s_3)(\cos A_1) \]
  to obtain $s_1 = 2.7724$\,rad.
  % FIGURE NEEDED: sphere with Lima, Yogyakarta, flight path and the
  % spherical triangle with sides s1, s2, s3 and angles A1, A2, A3
  % (2015_th_sol.pdf, page 20).
\item[2.] \emph{Determine the angles of the triangle} \hfill(30 points)\\
  Then use
  \[ \cos s_2 = (\cos s_1)(\cos s_3) + (\sin s_1)(\sin s_3)(\cos A_2) \]
  to obtain $A_2 = 0.35921$\,rad. Similarly, use the Law of Cosines
  \[ \cos s_3 = (\cos s_1)(\cos s_2) + (\sin s_1)(\sin s_2)(\cos A_3) \]
  to get that $A_3 = 0.3642$\,rad.
\item[3.] \emph{Use the Law of Cosines for sides to construct an equation in
  $\cos s_2'$.} \hfill(20 points)\\
  If P is the southernmost point, then the angle $A_3$ is a right angle,
  $A_3 = \frac{\pi}{2}$, the Law of
  \[ \cos s_3 = \cos s_1' \cos s_2' + \sin s_1' \sin s_2' \cos A_3
     = \cos s_1' \cos s_2' \]
  Thus,
  \[ \cos s_1' = \frac{\cos s_3}{\cos s_2'}
     = \frac{0.13538}{\cos s_2'}. \]
  Then substitute it to get
  \begin{align*}
  \cos s_2' &= \cos s_1' \cos s_3 + \sin s_1' \sin s_3 \cos A_2\\
    &= \frac{\cos s_3}{\cos s_2'}\cos s_3
     + \sqrt{1 - \left(\frac{\cos s_3}{\cos s_2'}\right)^{2}}\,
       \sin s_3 \cos A_2
  \end{align*}
\item[4.] \emph{Solve the equation} \hfill(20 points)\\
  Let $x = \cos s_2'$. Then
  \[ x = \frac{0.018328}{x}
     + \sqrt{\frac{x^2 - 0.018328}{x^2}} \times 0.92756 \]
  Multiply both sides by $x$ to get
  \begin{align*}
  x^2 &= 0.018328
     + x\sqrt{\frac{x^2 - 0.018328}{x^2}} \times 0.92756\\
  x^2 - 0.018328 &= x\sqrt{\frac{x^2 - 0.018328}{x^2}} \times 0.92756\\
  (x^2 - 0.018328)^2 &= (x^2 - 0.018328) \times (0.92756)^2\\
  x^4 - 0.89702x^2 + 0.016105 &= 0
  \end{align*}
  which is a quadratic equation of $x^2$. The roots are $\pm 0.13538$ and
  $\pm 0.93739$.
\item[5.] \emph{Determine the latitude} \hfill(10 points)\\
  Thus,
  \[ s_2' = \cos^{-1} 0.93739 = 0.35574\ \mathrm{rad} = 20^\circ 23' \]
  or
  \[ s_2' = \cos^{-1} 0.13538 = 1.4350\ \mathrm{rad} = 82^\circ 13' \]
  The latitude is $69^\circ 37'\,S$ or $7^\circ 47'$\,S. The second one is
  impossible because it is higher than the latitude of Lima. Thus the answer
  is
  \[ 69^\circ 37'\,S \]
\end{description}

\textbf{3rd version (using Napier's rule)}

\begin{description}
\item[1.] \emph{Determine the lengths of the sides of the triangle}
  \hfill(20 points)\\
  \[ s_2 = 90^\circ - 12^\circ 2' = 77^\circ 58' = 1.3608\ \mathrm{rad},\quad
     s_3 = 90^\circ - 7^\circ 47' = 82^\circ 13' = 1.4350\ \mathrm{rad}; \]
  \[ A_1 = (180^\circ - 77^\circ 1') + (180^\circ - 110^\circ 26')
     = 172^\circ 33' = 3.0116\ \mathrm{rad} \]
  Using the Law of Cosines for sides
  \[ \cos s_1 = (\cos s_2)(\cos s_3) + (\sin s_2)(\sin s_3)(\cos A_1) \]
  to obtain $s_1 = 2.7724$\,rad.
\item[2.] \emph{Determine the angles of the triangle} \hfill(30 points)\\
  Then use
  \[ \cos s_2 = (\cos s_1)(\cos s_3) + (\sin s_1)(\sin s_3)(\cos A_2) \]
  to obtain $A_2 = 0.35921$\,rad. Similarly, use the Law of Cosines
  \[ \cos s_3 = (\cos s_1)(\cos s_2) + (\sin s_1)(\sin s_2)(\cos A_3) \]
  to get that $A_3 = 0.3642$\,rad.
\item[3.] \emph{Find the latitude of the point} \hfill(50 points)\\
  Let $P$ be the southernmost point on the flight path. Since the triangle
  has a right angle at P, then use one of Napier's rules
  \[ \sin s_2' = \sin s_3 \sin A_2 \]
  to get $s_2' = 0.35576\ \mathrm{rad} = 20^\circ 23'$. Therefore, the
  latitude of $P$ is
  \[ 90^\circ - 20^\circ 23' = 69^\circ 37'\,S \]
  % FIGURE NEEDED: sphere with flight path, southernmost point P and the
  % right spherical triangle at P (2015_th_sol.pdf, page 22).
\end{description}
